% Lösungen Einheit 1 von 4 – Abstand zweier Punkte (zusammenbau v0.1)
\einheitenkopf{Einheit 1 von 4 · Abstand zweier Punkte}

\erg{9}{a) hin: einen Abstand berechnen}
\erg{10}{a) Punkt–Punkt: Betrag des Verbindungsvektors \quad b) $\overrightarrow{AB} = (2 | 2 | 1)$, $d = \sqrt{4 + 4 + 1} = 3$ \quad c) $|\overrightarrow{PQ}| = |(60 | 80 | 0)| = 100$ m, $v = 100 : 4 = 25$ m/s $= 90$ km/h \quad d) S liegt senkrecht über P: S(5 | 7 | 8); T liegt diametral gegenüber: $T = M' - (3 | 4 | 0)$ mit M'(2 | 3 | 8), also T(−1 | −1 | 8) \quad e) Katheten gleich: $|\overrightarrow{AB}| = 3$, $|\overrightarrow{AC}| = \sqrt{5}\,t$, also $t = \tfrac{3}{\sqrt{5}} \approx 1{,}34$ \quad f) Gleichschenklig im rechten Winkel: $|\overrightarrow{AB_s}| = |\overrightarrow{AC_t}|$, also $6s = 3t$ und $t = 2s$ – das ist die abgebildete Ursprungsgerade mit der Steigung $m = 2$ \quad g) $1 + 2a - a^2 = 0$ ergibt $a = 1 + \sqrt{2}$ (die negative Lösung liegt vor der Düse); W(7,24 | 9,66 | 0), Fußpunkt der Düse F(0 | 0 | 0); Weite $|\overrightarrow{FW}| = 5a \approx 12{,}1$ m \quad h) falsch: Ein Punkt von L mit $y = 5$ hat $z = 0$, also P(x | 5 | 0); $|\overrightarrow{HP}|^2 = (x-3)^2 + 16 + 16 \ge 32$, der Abstand zu H ist mindestens $\sqrt{32} \approx 5{,}66 > 5$}
\erg{11}{a) Fehler: Koordinaten addiert statt subtrahiert. Richtig: $\overrightarrow{AB} = (1 | 2 | 2)$, $|\overrightarrow{AB}| = \sqrt{1 + 4 + 4} = 3$ \quad b) Die Koordinatendifferenzen sind die Kanten eines achsenparallelen Quaders, die Verbindungsstrecke ist seine Raumdiagonale; zweimal Pythagoras (erst Flächen-, dann Raumdiagonale) ergibt $d^2 = a^2 + b^2 + c^2$. \quad c) $|\overrightarrow{TB}| = \sqrt{1\,200^2 + 900^2 + 1\,200^2} \approx 1\,921$ m; Fahrzeit $\approx 384$ s, also etwa 6,4 min}
